\newpage
\thispagestyle{fancy}

\begin{center}
  {\bfseries\large PRESENTACIÓN}
\end{center}

\vspace{0.15cm}

\begingroup
\justifying
\fontsize{9.5pt}{13.2pt}\selectfont

El presente documento constituye el estudio de factibilidad y formulación técnica conforme a la metodología oficial de Planifica Ecuador para el proyecto de inversión pública: \textit{Conectividad vial y construcción del puente vehicular sobre el río Santa Rosa en el sector El Guayabo - La Chilca}, en la parroquia rural Torata, cantón Santa Rosa, provincia de El Oro. Esta iniciativa surge como respuesta inaplazable a la incomunicación vial recurrente que enfrentan las familias campesinas y productores agropecuarios asentados en la cuenca alta del río Santa Rosa.

La problemática se asocia a dos factores estructurales: la ausencia definitiva de un puente formal vehicular con diseño hidráulico y gálibo normativo sobre el afluente, y el régimen hidrológico torrencial invernal que destruye anualmente los pasos rústicos de madera y sobrepasa el vado de piedra provisional. Esta situación provoca el aislamiento físico de las comunidades durante más de 105 días al año, generando el derrame o pérdida de más de 3.500 litros diarios de leche cruda, el incremento del 59 \% en el costo de los fletes de carga rural y la interrupción recurrente del acceso a los servicios esenciales de educación y salud.

El proyecto es impulsado mancomunadamente por el Gobierno Autónomo Descentralizado Provincial de El Oro en concurrencia con el GAD Municipal de Santa Rosa y el GAD Parroquial Rural de Torata, en articulación directa con las prioridades de fomento vial establecidas en el Plan de Desarrollo y Ordenamiento Territorial (PDOT) 2023--2027. La intervención contempla la cimentación profunda de dos estribos masivos de hormigón armado, obras de protección hidráulica con gaviones y la instalación de una superestructura modular metálica tipo Bailey de 30 metros de luz libre, garantizando la transitabilidad continua y segura los 365 días del año.

Por lo que la propuesta plantea una intervención integral orientada al desarrollo territorial mediante:
\begin{itemize}\setlength\itemsep{-0.2em}\setlength\parskip{0pt}
  \item La construcción e implantación de un puente modular tipo Bailey de 30 metros de longitud (AASHTO HS-20) con capacidad portante para transporte pesado lechero y ganadero.
  \item La ejecución de obras de estabilización de márgenes ribereñas mediante muros de gaviones y conformación técnica de los terraplenes de acceso en 150 metros a cada extremo del cauce.
  \item Un modelo de gestión y mantenimiento participativo articulado al Comité Vial Comunitario El Guayabo - La Chilca y al GAD Parroquial Rural de Torata.
\end{itemize}

El análisis de mercado desarrollado evidencia una demanda insatisfecha absoluta de 105 días de incomunicación total por año, donde la oferta vial formal en época invernal cae a cero. A partir del primer año de puesta en marcha, la infraestructura restituye el 100 \% de los días habilitados de paso seguro, erradicando los trasbordos a pie y posibilitando el flujo permanente de más de 3.650 litros diarios de leche y 60 vehículos por día hacia los mercados de Santa Rosa y Machala.

Finalmente, este proyecto integra análisis demográficos, técnicos, sociales, ambientales y de evaluación socioeconómica que demuestran su plena viabilidad y rentabilidad para el Estado ecuatoriano, garantizando una mejora sustancial en los ingresos autónomos y en la calidad de vida de las familias rurales de Torata.

\par
\endgroup

\clearpage
